\documentclass[11pt]{amsart}

\usepackage[margin=1.06in]{geometry}
\usepackage{amsmath,amssymb,amsthm}
\usepackage{enumitem}
\usepackage{graphicx}
\usepackage[hidelinks]{hyperref}

\newtheorem{theorem}{Theorem}
\newtheorem{lemma}{Lemma}
\theoremstyle{remark}
\newtheorem{remark}{Remark}

\newcommand{\calT}{\mathcal T}
\newcommand{\calL}{\mathcal L}
\newcommand{\conv}{\operatorname{conv}}
\newcommand{\wot}{\operatorname{WOT}}
\newcommand{\ip}[2]{\langle #1,#2\rangle}

\title{A block-norm counterexample in the $L^p(L^q)$ dilation problem}
\author{Agent lane 04}
\date{Run \texttt{fa\_banach\_001}, 2026-06-29}

\begin{document}
\maketitle

\section*{Status}

This packet gives a candidate counterexample to the equality question in
Open Problem 7.6 of Fackler--Gluck \cite{FacklerGlueck2017}.  Let
\[
        X_{p,q}=L^p([0,1];L^q([0,1])),\qquad 1<p,q<\infty,
\]
and let $\calT_{p,q}$ be the weak operator closed convex hull of the positive
invertible isometries of $X_{p,q}$.  The packet proves, using the standard
Guerre--Raynaud fibre form for positive isometries of $L^p(L^q)$, that
\[
        \calT_{p,q}\ne\{T\in\calL(X_{p,q}): T\ge 0,\ \|T\|\le 1\}
        \qquad (p\ne q).
\]
For $p=q$, Fubini identifies $X_{p,p}$ with scalar $L^p([0,1]^2)$, and the
scalar result used in the source gives equality.  Thus the first question in
Open Problem 7.6 has the sharp yes/no answer: equality holds for $p=q$ and
fails for $p\ne q$.  The characterization question for $p\ne q$ remains open.

\section*{Source open problem}

The source asks:
\begin{quote}
Let $p,q\in(1,\infty)$ and let $\mathcal T$ denote the weak operator closure
of the convex hull of all positive invertible isometries on
$L^p([0,1];L^q([0,1]))$.  Does $\mathcal T$ coincide with the set of all
positive contractions?  If not, can a good characterization of the elements of
$\mathcal T$ be given?
\end{quote}

\begin{center}
\includegraphics[width=0.94\textwidth]{figures/open_problem_crop.png}
\end{center}

\section*{Idea of the proof}

Operators in $\calT_{p,q}$ inherit a finite-dimensional shadow from the
fibre form of positive invertible isometries.  Compress such an operator to a
rectangular step sublattice
\[
        E=\ell_p^2(\ell_q^2)\subset L^p([0,1];L^q([0,1])).
\]
If $R\in\calT_{p,q}$ and $R_{ij}:\ell_q^2\to\ell_q^2$ are the four outer
blocks of this compression, then the scalar matrix
\[
        (\|R_{ij}\|)_{i,j=1}^2
\]
must be dominated by a positive contraction on $\ell_p^2$.  The reason is that
positive invertible isometries act as an outer scalar $L^p$ isometry together
with fibrewise positive $L^q$ isometries; the outer scalar compression supplies
the dominating $\ell_p$ contraction, and convexity plus weak closure preserve
the domination.

We then write down a positive contraction on $\ell_p^2(\ell_q^2)$ whose block
norm matrix is either
\[
        \begin{pmatrix}1&1\\0&0\end{pmatrix}
        \quad(q>p),
        \qquad\text{or}\qquad
        \begin{pmatrix}1&0\\1&0\end{pmatrix}
        \quad(p>q).
\]
These matrices have $\ell_p^2$ operator norm $2^{1/p'}$ and $2^{1/p}$,
respectively, both strictly larger than one.  The lifted positive contraction
therefore cannot belong to $\calT_{p,q}$.

\section*{The finite compression obstruction}

Let $A_1,A_2$ be a measurable partition of $[0,1]$ into two sets of positive
measure, and let $B_1,B_2$ be another such partition.  Let
$E=\ell_p^2(\ell_q^2)$, with vectors written
\[
        x=((x_{11},x_{12}),(x_{21},x_{22})).
\]
Embed $E$ isometrically into $X_{p,q}$ as the rectangular step sublattice by
using normalized indicators of the four rectangles $A_i\times B_\alpha$.
Denote this embedding by $J:E\to X_{p,q}$.  Let $P$ be conditional expectation
onto this rectangular sigma algebra and set $Q=J^{-1}P$.  Then $J$ is a
positive isometry, $Q$ is a positive contraction, and $QJ=I_E$.

\begin{lemma}[block domination]
\label{lem:block-domination}
Assume $p\ne q$.  If $R\in\calT_{p,q}$ and
\[
        QRJ=(R_{ij})_{i,j=1}^2
\]
is written as a $2\times2$ matrix of operators
$R_{ij}\in\calL(\ell_q^2)$, then the scalar matrix
$B_R=(\|R_{ij}\|)_{i,j=1}^2$ satisfies
\[
        \|B_R\|_{\ell_p^2\to\ell_p^2}\le 1.
\]
\end{lemma}

\begin{proof}
It suffices first to prove the assertion for a positive invertible isometry
$U$ of $X_{p,q}$.  We use the standard Guerre--Raynaud form for positive
isometries of $L^p(L^q)$, in the form quoted for $L_p(L_q)$ in
\cite[proof of Proposition 6.12]{CuthDeRancourtDoucha2024}: a positive
invertible isometry acts by a positive scalar $L^p$ isometry on the outer
variable and by a measurable field of positive $L^q$ isometries on the fibres.
Equivalently, there are a positive invertible isometry $S$ of scalar
$L^p[0,1]$ and fibre isometries $V_s$ such that the outer norm function of
$Uf$ is the scalar isometric image under $S$ of the outer norm function of $f$.

Let $J_p:\ell_p^2\to L^p[0,1]$ be the normalized two-step embedding associated
with $A_1,A_2$, and let $Q_p:L^p[0,1]\to\ell_p^2$ be the corresponding
conditional expectation quotient.  The scalar compression
\[
        C=Q_pSJ_p
\]
is a positive contraction on $\ell_p^2$; write $C=(c_{ij})_{i,j=1}^2$.
For $x\in\ell_q^2$ placed in the $j$th outer coordinate of $E$, the fibre
isometries do not increase the subsequent conditional expectation onto
$\operatorname{span}\{\chi_{B_1},\chi_{B_2}\}$.  Hence the $i,j$ block of
$QUJ$ has norm at most $c_{ij}$:
\[
        \|(QUJ)_{ij}x\|_{\ell_q^2}\le c_{ij}\|x\|_{\ell_q^2}.
\]
Thus the block norm matrix of $QUJ$ is dominated entrywise by the positive
$\ell_p^2$ contraction $C$.

Now let $R$ be in the convex hull of positive invertible isometries.  The
block norm matrix of $QRJ$ is entrywise dominated by the same convex
combination of the corresponding scalar compression matrices, hence by a
positive $\ell_p^2$ contraction.  Finally, if $R_\alpha\to R$ weakly in
$\calL(X_{p,q})$, then $QR_\alpha J\to QRJ$ weakly on the finite-dimensional
space $E$, hence in norm.  The block norms therefore converge, so the
conclusion passes to the weak operator closure.
\end{proof}

\section*{The counterexample}

\begin{theorem}
\label{thm:counterexample}
Let $1<p,q<\infty$ and $p\ne q$.  Then $\calT_{p,q}$ is a proper subset of
the positive contractions on $X_{p,q}$.
\end{theorem}

\begin{proof}
We construct the finite model first, then lift it through $J$ and $Q$.

Suppose first that $q>p$.  Define $S_>:\ell_p^2(\ell_q^2)\to
\ell_p^2(\ell_q^2)$ by
\[
        S_>((a,b),(c,d))=((a,d),(0,0)).
\]
This is positive.  Moreover
\[
 \|S_>((a,b),(c,d))\|
 = (|a|^q+|d|^q)^{1/q}
 \le \big(\|(a,b)\|_q^p+\|(c,d)\|_q^p\big)^{1/p},
\]
because $q>p$ implies $\|\cdot\|_{\ell_q^2}\le\|\cdot\|_{\ell_p^2}$ on
non-negative two-vectors.  Hence $S_>$ is a positive contraction.  Its outer
blocks have norm matrix
\[
        B_{S_>}=\begin{pmatrix}1&1\\0&0\end{pmatrix},
\]
whose norm on $\ell_p^2$ is $2^{1/p'}>1$.

Suppose next that $p>q$.  Define $S_<:\ell_p^2(\ell_q^2)\to
\ell_p^2(\ell_q^2)$ by
\[
        S_<((a,b),(c,d))=((a,0),(0,b)).
\]
Again $S_<$ is positive, and
\[
 \|S_<((a,b),(c,d))\|
 = (|a|^p+|b|^p)^{1/p}
 \le (|a|^q+|b|^q)^{1/q}
 \le \big(\|(a,b)\|_q^p+\|(c,d)\|_q^p\big)^{1/p},
\]
so $S_<$ is a positive contraction.  Its block norm matrix is
\[
        B_{S_<}=\begin{pmatrix}1&0\\1&0\end{pmatrix},
\]
whose norm on $\ell_p^2$ is $2^{1/p}>1$.

In either case call the relevant finite contraction $S$.  Set
\[
        T=JSQ\in\calL(X_{p,q}).
\]
Since $J$ is positive isometric, $Q$ is a positive contraction, and $S$ is a
positive contraction, $T$ is a positive contraction on $X_{p,q}$.  But
\[
        QTJ=Q(JSQ)J=S.
\]
If $T$ belonged to $\calT_{p,q}$, Lemma~\ref{lem:block-domination} would imply
that the block norm matrix of $S$ has $\ell_p^2$ norm at most one.  The
explicit matrices above have norms $2^{1/p'}$ or $2^{1/p}$, a contradiction.
Thus $T\notin\calT_{p,q}$.
\end{proof}

\section*{Consequences and limitations}

For $p=q$, Fubini identifies $X_{p,p}$ with scalar $L^p([0,1]^2)$, and the
source recalls the scalar Greim theorem that positive contractions on atomless
scalar $L^p$ are weak operator limits of convex combinations of positive
invertible isometries.  Thus equality holds in the diagonal case $p=q$.
Theorem~\ref{thm:counterexample} shows equality fails off the diagonal.

This does not provide a full characterization of $\calT_{p,q}$ when $p\ne q$.
The obstruction suggests that such a characterization must include
outer-block domination constraints, not merely positivity and contractivity.

\section*{Novelty and verification notes}

\begin{itemize}[leftmargin=*]
\item Local indexes were searched for \texttt{1709.08547},
\texttt{L\^{}p(L\^{}q)}, positive invertible isometries, and counterexample
phrases.  The only exact prior packet was the earlier tensor-product partial
result for the same source problem.
\item Bounded web search on 2026-06-29 for the source title, Open Problem 7.6,
positive invertible isometries on $L^p(L^q)$, and weak operator closure found
no later paper giving this counterexample or a full characterization.
\item The finite norm checks are explicit: the two displayed block matrices
have norms $2^{1/p'}$ and $2^{1/p}$ on $\ell_p^2$.
\item The main human-review point is Lemma~\ref{lem:block-domination},
specifically the use of the standard Guerre--Raynaud fibre representation for
positive invertible isometries of $L^p(L^q)$ throughout the stated exponent
range $p\ne q$.
\end{itemize}

\section*{Human review recommendation}

Review as a likely valid counterexample and negative answer to the equality
part of Open Problem 7.6.  If the Guerre--Raynaud fibre form is accepted in
the stated positive-isometry range, the rest of the proof is a finite
two-by-two norm computation.

\begin{thebibliography}{9}

\bibitem{FacklerGlueck2017}
S. Fackler and J. Gluck,
\emph{A Toolkit for Constructing Dilations on Banach Spaces},
arXiv:1709.08547, 2017.
\url{https://arxiv.org/abs/1709.08547}

\bibitem{CuthDeRancourtDoucha2024}
M. C\'uth, N. de Rancourt, and M. Doucha,
\emph{Guarded Fraisse Banach spaces},
arXiv:2408.12755, 2024.
\url{https://arxiv.org/abs/2408.12755}

\bibitem{GuerreRaynaud1988}
S. Guerre and Y. Raynaud,
\emph{Sur les isometries de \(L^p(X)\) et le theoreme ergodique vectoriel},
Canad. J. Math. 40 (1988), no. 2, 360--391.

\end{thebibliography}

\end{document}
